\documentclass[11pt]{article}

\usepackage[margin=1in]{geometry}
\usepackage{amsmath,amssymb}
\usepackage{booktabs}
\usepackage{array}
\usepackage{enumitem}
\usepackage[hidelinks]{hyperref}

\newcommand{\expit}{\operatorname{expit}}
\newcommand{\ind}{\mathbb{I}}
\newcommand{\Prb}{\operatorname{Pr}}
\newcommand{\RMST}{\operatorname{RMST}}

\title{Combined Site Heterogeneity\\
Simulation \& Method Note}
\author{}
\date{}

\begin{document}
\maketitle

\section{Purpose}

This simulation evaluates federated clone--censor--weight (Fed-CCW) estimation
under three sources of between-site heterogeneity occurring simultaneously:
\begin{enumerate}[leftmargin=2em]
  \item unbalanced site sample sizes;
  \item different site-specific distributions of measured baseline
        covariates;
  \item different site treatment-initiation practices after
        conditioning on measured patient history.
\end{enumerate}
Each low, moderate, or high scenario specifies all three components jointly.
The total sample size remains 5,000, but site sizes become increasingly unequal. 
Patient-mix separation and initiation practice separation increase at the same time.

The simulation mainly compares Fed-CCW with two pooled CCW implementations (along with some additional naive methods). 
The first pooled implementation uses fully site-stratified nuisance models and is thus mathematically aligned with Fed-CCW. 
The second uses a conventional pooled logistic initiation model with site fixed effects and one common covariate slope vector, thereby borrowing information across sites.


\section{Indices, observed data, and simulation grid}

Sites are indexed by $k\in\{1,\ldots,5\}$, patients by $i$, and discrete
follow-up intervals by $m=1,\ldots,M$, where $M=25$. For patient $i$ at site
$k$, let
\begin{itemize}[leftmargin=2em]
  \item $X_{1ik}$ and $X_{2ik}$ denote baseline covariates;
  \item $L_{1ikm}$ and $L_{2ikm}$ denote time-varying covariates;
  \item $S_{ik}$ denote the treatment-initiation interval;
  \item $T_{ik}$ denote the terminal-event interval; and
  \item $A_{ik}^{\tau}=\ind(S_{ik}\leq\tau)$ indicate initiation within the
        grace period.
\end{itemize}
If initiation or the event does not occur during follow-up, its time is coded
as $M+1$.

The simulation grid is
\[
  \tau\in\{5,10,15\},\qquad
  \beta_{\mathrm{trt}}\in\{-1.0,-0.7,-0.5\},\qquad
  h\in\{\text{low},\text{moderate},\text{high}\},
\]
giving $3\times3\times3=27$ array tasks. 

\section{Joint site-heterogeneity settings}

For heterogeneity level $h$, define
\[
  \mathcal G_h
  =\left\{n_{kh},\mu_{kh},\sigma_{kh},p_{kh},\alpha_{kh}:
            k=1,\ldots,5\right\},
\]
where $n_{kh}$ is site size, $X_1\sim N(\mu_{kh},\sigma_{kh}^2)$,
$X_2\sim\operatorname{Bernoulli}(p_{kh})$, and $\alpha_{kh}$ is the
site-specific treatment-initiation intercept.

\begin{table}[ht]
\centering
\small
\caption{Combined site-heterogeneity settings.}
\begin{tabular}{llrrrr}
\toprule
Level & Site & $n_{kh}$ & $X_1$ distribution & $X_2$ probability & $\alpha_{kh}$ \\
\midrule
Low & 1 & 900  & $N(-0.250,0.950^2)$ & 0.350 & $-3.250$ \\
    & 2 & 950  & $N(-0.125,0.975^2)$ & 0.375 & $-3.125$ \\
    & 3 & 1000 & $N(0,1^2)$          & 0.400 & $-3.000$ \\
    & 4 & 1050 & $N(0.125,1.025^2)$  & 0.425 & $-2.875$ \\
    & 5 & 1100 & $N(0.250,1.050^2)$  & 0.450 & $-2.750$ \\
\addlinespace
Moderate & 1 & 600  & $N(-0.750,0.850^2)$ & 0.200 & $-3.750$ \\
         & 2 & 800  & $N(-0.375,0.925^2)$ & 0.300 & $-3.375$ \\
         & 3 & 1000 & $N(0,1^2)$          & 0.400 & $-3.000$ \\
         & 4 & 1200 & $N(0.375,1.075^2)$  & 0.500 & $-2.625$ \\
         & 5 & 1400 & $N(0.750,1.150^2)$  & 0.600 & $-2.250$ \\
\addlinespace
High & 1 & 400  & $N(-1.500,0.700^2)$ & 0.100 & $-4.500$ \\
     & 2 & 600  & $N(-0.750,0.850^2)$ & 0.250 & $-3.750$ \\
     & 3 & 800  & $N(0,1^2)$          & 0.400 & $-3.000$ \\
     & 4 & 1200 & $N(0.750,1.150^2)$  & 0.550 & $-2.250$ \\
     & 5 & 2000 & $N(1.500,1.300^2)$  & 0.700 & $-1.500$ \\
\bottomrule
\end{tabular}
\end{table}

For every level,
\[
  \sum_{k=1}^5 n_{kh}=5000.
\]
The maximum-to-minimum site-size ratios are approximately 1.22, 2.33, and
5.00 for the low, moderate, and high levels. The $X_1$ mean ranges are 0.5,
1.5, and 3.0; the $X_2$ probability ranges are 0.1, 0.4, and 0.6; and the
initiation-intercept ranges are 0.5, 1.5, and 3.0.

The gradients are intentionally aligned. Larger sites tend to enroll patients
with larger $X$ values and have higher residual initiation propensity. 
% Thus, high heterogeneity is a joint stress test rather than three independent one-factor perturbations.

\section{Data-generating process}

\subsection{Baseline covariates}

At level $h$, site $k$ generates
\[
  X_{1ik}\sim N(\mu_{kh},\sigma_{kh}^2),
  \qquad
  X_{2ik}\sim\operatorname{Bernoulli}(p_{kh}).
\]
Because $n_{kh}$ also varies, the marginal target-population patient mix is
the size-weighted five-site mixture rather than an equally weighted mixture.

\subsection{Time-varying covariates}

For $\ell\in\{1,2\}$,
\begin{equation}
\label{eq:L-dgp}
  L_{\ell,ikm}
  =0.6L_{\ell,ik,m-1}
   +0.3X_{1ik}
   +0.2X_{2ik}
   +\varepsilon_{\ell,ikm},
  \qquad
  \varepsilon_{\ell,ikm}\sim N(0,1),
\end{equation}
with $L_{\ell,ik0}=0$. Treatment does not affect subsequent $L$. Covariates
are generated only while the patient remains alive.

\subsection{Treatment initiation}

Among patients alive and untreated at the beginning of interval $m$, define
\[
  q_{ikm}
  =\Prb(S_{ik}=m\mid S_{ik}\geq m,T_{ik}\geq m,
                         X_{ik},L_{ikm},K_i=k).
\]
The initiation DGP is
\begin{equation}
\label{eq:init-dgp}
  \operatorname{logit}(q_{ikm})
  =\alpha_{kh}
   +0.8X_{1ik}
   -0.3X_{2ik}
   +0.5L_{1ikm}
   +0.4L_{2ikm}.
\end{equation}
The covariate slopes are common across sites and heterogeneity levels. Only
the site intercept $\alpha_{kh}$ changes. Consequently, two patients with the
same measured history can have different initiation probabilities because
they receive care at different sites.

\subsection{Terminal-event process}

Conditional on being alive at interval $m$, the event hazard is
\begin{align}
\label{eq:event-dgp}
  \operatorname{logit}
  \Prb(T_{ik}=m\mid T_{ik}\geq m,\mathcal H_{ikm})
  ={}&-2.5+0.5X_{1ik}+0.4X_{2ik} \notag\\
     &+0.4L_{1ikm}+0.3L_{2ikm}
      +0.2L_{1ik,m-1}+0.15L_{2ik,m-1} \notag\\
     &+\beta_{\mathrm{trt}}
       \ind(S_{ik}\leq\tau,m\geq S_{ik}).
\end{align}
Treatment affects the event hazard only if initiation occurs within the grace period. 
There is no additional direct site term in the event model. 
Site heterogeneity influences the marginal outcome distribution through patient mix and treatment initiation, 
not through an independent site outcome effect (so no treatment heterogeneity induced by site).

\medskip

Within each interval, the simulation generates current $L$, then treatment
initiation, and then the event. A patient may initiate and experience the event
in the same interval. Initiation and covariate evolution stop after death.

\section{Strategies, cloning, and artificial censoring}

The two strategies are
\begin{align*}
  g=1:&\quad \text{initiate treatment within the grace period},\\
  g=0:&\quad \text{do not initiate treatment within the grace period}.
\end{align*}
For the $g=1$ clone,
\[
  C^1_{ik}
  =\begin{cases}
      M+1, & S_{ik}\leq\tau,\\
      \tau-1, & S_{ik}>\tau.
    \end{cases}
\]
For the $g=0$ clone,
\[
  C^0_{ik}
  =\begin{cases}
      S_{ik}-1, & S_{ik}\leq\tau,\\
      M+1, & S_{ik}>\tau.
    \end{cases}
\]
Deviation occurs during the interval that triggers artificial censoring, so
the clone is retained only through the preceding interval.

\section{Oracle truth}

\subsection{Oracle target population}

For each $(\tau,\beta_{\mathrm{trt}},h)$ cell, the oracle simulates
$N_{\mathrm{oracle}}=3{,}000{,}000$ observations. Its site proportions match
the finite-sample target population:
\begin{equation}
\label{eq:site-weight}
  \pi_{kh}=\frac{n_{kh}}{\sum_{j=1}^5 n_{jh}}
           =\frac{n_{kh}}{5000}.
\end{equation}
Therefore, the oracle changes with the heterogeneity level not only because
the site-specific DGP changes, but also because the target-population site
mixture changes.

\subsection{Oracle denominator and common numerator}

The oracle uses the true initiation probability
\[
  q_{ikm}^{D,\mathrm{true}}
  =\expit\!\left(
      \alpha_{kh}+0.8X_{1ik}-0.3X_{2ik}
      +0.5L_{1ikm}+0.4L_{2ikm}
    \right).
\]
Its common marginal numerator hazard is
\begin{equation}
\label{eq:oracle-num}
  q_m^{N,\mathrm{oracle}}
  =\frac{
      \sum_{k=1}^5\sum_i
      \ind(S_{ik}=m,T_{ik}\geq m)
    }{
      \sum_{k=1}^5\sum_i
      \ind(S_{ik}\geq m,T_{ik}\geq m)
    }.
\end{equation}
This numerator automatically reflects unequal site sizes, different patient
mixes, and different initiation practices.

\subsection{Oracle weights and survival}

Writing $q_m^N=q_m^{N,\mathrm{oracle}}$ and
$q_{ikm}^D=q_{ikm}^{D,\mathrm{true}}$, the waiting and initiation factors are
\[
  r^{\mathrm{wait}}_{ikm}
  =\frac{1-q_m^N}{1-q_{ikm}^D},
  \qquad
  r^{\mathrm{init}}_{ikm}
  =\frac{q_m^N}{q_{ikm}^D}.
\]
The initiate-within-$\tau$ weight is
\begin{equation}
\label{eq:w1}
  W^1_{ikm}
  =\prod_{j=1}^{\min(m,\tau)}
    \left\{
      \left(\frac{q_j^N}{q_{ikj}^D}\right)^{\ind(S_{ik}=j)}
      \left(\frac{1-q_j^N}{1-q_{ikj}^D}\right)^{\ind(S_{ik}>j)}
    \right\},
\end{equation}
while the do-not-initiate weight is
\begin{equation}
\label{eq:w0}
  W^0_{ikm}
  =\prod_{j=1}^{\min(m,\tau,S_{ik}-1)}
     \frac{1-q_j^N}{1-q_{ikj}^D}.
\end{equation}
Factors are evaluated only while the patient remains eligible. After $\tau$,
the weights are carried forward. The current setting uses no effective weight
truncation.

For strategy $g$, define
\[
  R^g_{ikm}=\ind(T_{ik}>m-1,C^g_{ik}>m-1)
\]
and
\[
  D^g_{ikm}=\ind(R^g_{ikm}=1,T_{ik}=m,T_{ik}\leq C^g_{ik}).
\]
The oracle hazard and survival curve are
\[
  \lambda_m^{g,\mathrm{oracle}}
  =\frac{\sum_{k,i}W^g_{ikm}D^g_{ikm}}
         {\sum_{k,i}W^g_{ikm}R^g_{ikm}},
  \qquad
  S^{g,\mathrm{oracle}}(m)
  =\prod_{j=1}^m(1-\lambda_j^{g,\mathrm{oracle}}).
\]

At $M=25$,
\[
  \psi_g=1-S^g(M),
  \qquad
  \RMST_g=\sum_{m=0}^{M-1}S^g(m).
\]
The reported contrasts are
\[
  \mathrm{RD}=\psi_1-\psi_0,
  \qquad
  \mathrm{RR}=\frac{\psi_1}{\psi_0},
\]
\[
  \mathrm{OR}
  =\frac{\psi_1/(1-\psi_1)}{\psi_0/(1-\psi_0)},
  \qquad
  \RMST_{\mathrm{diff}}=\RMST_1-\RMST_0.
\]

\section{Federated CCW}

\subsection{Local denominator nuisance models}

Each site constructs its own eligible person-interval dataset and fits
\begin{align}
\label{eq:local-den}
  \operatorname{logit}(\widehat q^{D,\mathrm{local}}_{ikm})
  ={}&\widehat\gamma_{km}
      +\widehat\beta_{1k}X_{1ik}
      +\widehat\beta_{2k}X_{2ik}
      +\widehat\beta_{3k}L_{1ikm}
      +\widehat\beta_{4k}L_{2ikm} \notag\\
     &+\widehat\beta_{5k}L_{1ik,m-1}
      +\widehat\beta_{6k}L_{2ik,m-1}.
\end{align}
The interval effect $\widehat\gamma_{km}$ is represented by a factor for
$m$. Every coefficient is estimated separately within site $k$. Thus,
Fed-CCW does not borrow nuisance regression coefficients across sites. This
model is more flexible than the DGP because it permits site-specific slopes
and includes lagged $L$ in the initiation model.

\subsection{Common numerator and central estimation}

Site $k$ releases
\[
  N^{\mathrm{init}}_{km}
  =\sum_i\ind(S_{ik}=m,T_{ik}\geq m),
  \qquad
  N^{\mathrm{eligible}}_{km}
  =\sum_i\ind(S_{ik}\geq m,T_{ik}\geq m).
\]
The center computes and broadcasts
\begin{equation}
\label{eq:fed-num}
  \widehat q_m^N
  =\frac{\sum_kN^{\mathrm{init}}_{km}}
         {\sum_kN^{\mathrm{eligible}}_{km}}.
\end{equation}
Each site then constructs the strategy weights using its local denominator and
the common numerator. It releases weighted event and risk totals
\[
  DW^g_{km}=\sum_iW^g_{ikm}D^g_{ikm},
  \qquad
  RW^g_{km}=\sum_iW^g_{ikm}R^g_{ikm}.
\]
The central hazard is
\begin{equation}
\label{eq:fed-hazard}
  \widehat\lambda_m^{g,\mathrm{fed}}
  =\frac{\sum_kDW^g_{km}}{\sum_kRW^g_{km}},
\end{equation}
followed by
\[
  \widehat S^{g,\mathrm{fed}}(m)
  =\prod_{j=1}^m(1-\widehat\lambda_j^{g,\mathrm{fed}}).
\]
This gives one pooled weighted survival curve, and thus RMST, RD, RR, and OR estimates, for each strategy. 

For uncertainty estimation, global hazards and risk totals are returned to
the sites. Sites calculate local individual influence contributions and
release sums of squares. The implemented model-based variance treats the
estimated weights as fixed and does not add nuisance-model estimation
uncertainty.

\section{Pooled CCW with site-stratified nuisance models}

The aligned pooled estimator has all individual records centrally but fits
Equation~\eqref{eq:local-den} separately within each site. It uses the same
common numerator, artificial censoring, weight construction, weighted survival
estimator, and influence-function calculation as Fed-CCW.

Its pooled weighted hazard is
\[
  \widehat\lambda_m^{g,\mathrm{strat}}
  =\frac{\sum_k\sum_iW^g_{ikm}D^g_{ikm}}
         {\sum_k\sum_iW^g_{ikm}R^g_{ikm}}.
\]
Since the local sufficient statistics satisfy
\[
  DW^g_{km}=\sum_iW^g_{ikm}D^g_{ikm},
  \qquad
  RW^g_{km}=\sum_iW^g_{ikm}R^g_{ikm},
\]
we have
\[
  \widehat\lambda_m^{g,\mathrm{strat}}
  =\widehat\lambda_m^{g,\mathrm{fed}}.
\]
Therefore, Fed-CCW and site-stratified pooled CCW have identical survival
curves, treatment-effect estimates, and currently implemented model-based
standard errors up to numerical precision. This comparison verifies lossless
federation for the fully local nuisance-model specification.

\section{Pooled CCW with site fixed effects}

The second, more conventional, pooled estimator fits one logistic initiation model to all pooled
person-interval records:
\begin{align}
\label{eq:pooled-fe}
  \operatorname{logit}(\widehat q^{D,\mathrm{FE}}_{ikm})
  ={}&\widehat\gamma_m+\widehat\delta_k
      +\widehat\beta_1X_{1ik}+\widehat\beta_2X_{2ik} \notag\\
     &+\widehat\beta_3L_{1ikm}+\widehat\beta_4L_{2ikm}
      +\widehat\beta_5L_{1ik,m-1}
      +\widehat\beta_6L_{2ik,m-1}.
\end{align}
Here, $\widehat\gamma_m$ is a common interval effect,
$\widehat\delta_k$ is a site fixed effect, and the $\widehat\beta$ vector is
shared across sites. The model therefore permits residual differences in site
initiation propensity while borrowing information across sites for the
effects of patient history.

The common-slope restriction in Equation~\eqref{eq:pooled-fe} is correct under
the DGP, because Equation~\eqref{eq:init-dgp} uses common covariate
slopes and site-specific intercepts. The pooled fixed-effect model may
therefore produce more stable nuisance estimates and weights, particularly at
small sites. However, it is not the same statistical estimator as Fed-CCW's
fully local nuisance-model approach.

After fitting Equation~\eqref{eq:pooled-fe}, this estimator uses the same
common numerator and clone-censor-weight construction. Its hazard is
\[
  \widehat\lambda_m^{g,\mathrm{FE}}
  =\frac{\sum_{k,i}W^{g,\mathrm{FE}}_{ikm}D^g_{ikm}}
         {\sum_{k,i}W^{g,\mathrm{FE}}_{ikm}R^g_{ikm}}.
\]
Because $W^{g,\mathrm{FE}}_{ikm}$ generally differs from the locally estimated
Fed-CCW weight, this estimator need not equal Fed-CCW in finite samples. Any
efficiency advantage reflects the additional common-slope assumption and
cross-site information borrowing, not merely centralized data access.

\section{Other comparison methods}

The simulation also includes:
\begin{enumerate}[leftmargin=2em]
  \item federated IPW without cloning;
  \item federated unweighted per-protocol analysis; and
  \item federated landmark IPW.
\end{enumerate}
Together with the three CCW implementations, this gives six methods and eight
reported estimands per method: $\psi_1$, $\psi_0$, RD, RR, OR, $\RMST_1$,
$\RMST_0$, and $\RMST_{\mathrm{diff}}$.



\end{document}
